\section*{TÓM TẮT ĐỒ ÁN}
\phantomsection\addcontentsline{toc}{section}{\numberline {}TÓM TẮT ĐỒ ÁN}

Bài toán phát hiện gian lận thẻ tín dụng là một lĩnh vực quan trọng trong tài chính, nhằm xác định các hành vi bất thường ẩn trong tập dữ liệu lớn và mất cân bằng, nơi mà các giao dịch hợp lệ chiếm tỷ lệ áp đảo. Đặc trưng của bài toán này là dữ liệu có tính chất tuần tự và phi cấu trúc, đòi hỏi mô hình không chỉ trích xuất được các đặc trưng quan trọng mà còn phải nắm bắt được các mối quan hệ hành vi theo chuỗi thời gian.

Đồ án tốt nghiệp này tập trung nghiên cứu và triển khai một mô hình phát hiện gian lận thẻ tín dụng hiệu quả, không sử dụng kiến trúc đồ thị, mà thay vào đó là sự kết hợp giữa cơ chế chọn lọc đặc trưng (mask channel-wise) và mạng nơ-ron hồi tiếp (Recurrent Neural Network – RNN). Mô hình cũng tích hợp chiến lược mở rộng giao dịch (Transactional Expansion) để cung cấp thêm ngữ cảnh cho mỗi giao dịch, và sử dụng hàm mất mát Focal Loss để cải thiện khả năng học trong điều kiện dữ liệu mất cân bằng nghiêm trọng.

Kết quả thực nghiệm trên tập dữ liệu giả lập (tổng hợp) từ Kaggle cho thấy mô hình đề xuất đạt hiệu quả vượt trội so với các phương pháp truyền thống, với độ chính xác và F1-score cao. %Đồ án không chỉ trình bày cơ sở lý thuyết và kiến trúc mô hình, mà còn triển khai thực tế một hệ thống có khả năng ứng dụng trong các giải pháp giám sát gian lận tài chính quy mô lớn.


\newpage

\section*{THESIS SUMMARY}
\phantomsection\addcontentsline{toc}{section}{\numberline {}THESIS SUMMARY}

The problem of credit card fraud detection is a critical issue in the financial domain, aiming to identify fraudulent behavior hidden within large and imbalanced datasets, where legitimate transactions overwhelmingly dominate. This task is characterized by sequential and unstructured data, which requires models to capture complex relationships between transactions as well as user behavior patterns over time.

This graduation thesis focuses on developing an effective credit card fraud detection model that integrates feature selection mechanisms (via channel-wise masks) and recurrent neural networks (RNNs) to model temporal dependencies among transactions. Additionally, the model applies Transactional Expansion to enrich transaction context and Focal Loss to address the severe class imbalance issue commonly found in fraud detection.

Experimental results on the simulated (synthetic) dataset from Kaggle demonstrate that the proposed model outperforms traditional methods, achieving high accuracy and F1-score. %This work not only presents the theoretical foundation and model architecture but also implements a practical system that can be deployed in large-scale financial fraud monitoring environments.